\begin{exercise}[Self-financing wealth and the risk-neutral measure]
\textbf{Restatement.}
In the Black--Scholes model, express the discounted value of a self-financing
strategy \((x,\theta)\), compute its first two moments, and show that the
Girsanov measure with market price of risk \((\mu-r)/\sigma\) is risk-neutral.

\begin{solution}
We use the convention \(S^0_0=1\), so that initial capital \(x\) is also
initial discounted capital.  By self-financing, \polyref{Eq.~(3.8)} gives
\[
  \Vtilde_t^{x,\theta}
    =x+\int_0^t\theta_u\,\dd\Stilde_u
    =x+\int_0^t(\mu-r)\theta_u\Stilde_u\,\dd u
      +\int_0^t\sigma\theta_u\Stilde_u\,\dd B_u.
\]
If \(S^0_0\ne1\), replace \(x\) by \(x/S^0_0\) in this formula.

Set
\[
  A_t=(\mu-r)\int_0^t\theta_u\Stilde_u\,\dd u,
  \qquad
  M_t=\sigma\int_0^t\theta_u\Stilde_u\,\dd B_u.
\]
Under the usual square-integrability assumptions, \(M\) is a martingale with
\(\E[M_t]=0\), hence
\[
  \E[\Vtilde_t^{x,\theta}]
    =x+\E[A_t]
    =x+\E\left[\int_0^t(\mu-r)\theta_u\Stilde_u\,\dd u\right].
\]
The variance is
\[
  \Var(\Vtilde_t^{x,\theta})
    =\Var(A_t+M_t)
    =\Var(A_t)+\E\left[\int_0^t\sigma^2\theta_u^2\Stilde_u^2\,\dd u\right]
      +2\,\Cov(A_t,M_t).
\]
The covariance term is kept in the general adapted case; it vanishes only
under extra assumptions, for instance for deterministic integrands.

Now put
\[
  \lambda=\frac{\mu-r}{\sigma},
  \qquad
  Z_t=\exp\left(-\lambda B_t-\frac12\lambda^2t\right).
\]
The density must contain \(\lambda^2\) in the deterministic term.  Since
\(\lambda\) is constant, Novikov's condition is immediate and
\(\dd\Q/\dd\P|_{\F_t}=Z_t\) defines a probability measure equivalent to
\(\P\).  By Girsanov's theorem, stated in the poly before the BMS chapter and
used in \polyref{Eq.~(3.6)}, the process
\[
  B_t^\Q=B_t+\lambda t
\]
is a Brownian motion under \(\Q\).  Substituting
\(\dd B_t=\dd B_t^\Q-\lambda\dd t\) in
\(\dd S_t=\mu S_t\dd t+\sigma S_t\dd B_t\) gives
\[
  \dd S_t
    =(\mu-\sigma\lambda)S_t\dd t+\sigma S_t\dd B_t^\Q
    =rS_t\dd t+\sigma S_t\dd B_t^\Q.
\]
Therefore
\[
  \dd\Stilde_t=\sigma\Stilde_t\dd B_t^\Q,
\]
so the discounted risky asset is a \(\Q\)-martingale.  Thus \(\Q\) is a
risk-neutral probability measure.
\end{solution}
\end{exercise}
